% 实验评估
\providecommand{\todo}[1]{\textbf{[待真实数据：\ #1]}}% 实验数值占位符，待真实数据回填后删除

\chapter{实验评估}\label{chap:experiments}

第~\ref{chap:overview}章给出了统一 LP 框架（固定 $B$ 的可行性判定 + 二分 $B^*$），
第~\ref{chap:sensitivity}章给出了 $B^*$ 对期待与约束的灵敏度。
本章通过五个递进的实验问题验证框架的实用性：
框架能跑多大（可扩展性）、瓶颈在哪（$B^*$ 分析）、设计空间长什么样（Pareto 前沿）、
以及低精度判定何时足够（保真度）。
每个问题对应一个明确的论证目标，实验设计服务于论证逻辑。

需要说明：本章各实验的具体数值尚在获取中，因此本章只写清\textbf{实验设计}、
\textbf{论证逻辑}与\textbf{预期方向}，所有具体数值一律以 \todo{数值} 的形式占位标注，
待真实数据回填后替换。任何占位都不表示已被验证的结论，仅表示待验证的方向。

% 实验设置
\section{实验设置}\label{sec:setup}

硬件：Intel Core i7-13700K, 32GB RAM, Linux。
软件：Python 3.13, cvxpy 1.3 + CLARABEL求解器。
默认物理参数：12×12mm裸片，50W TDP功耗，45$\mu$m pitch $\mu$bump（75mA/bump），0.8V供电，
70\%面积利用率，液冷（2.0 W/mm$^2$），UCIe-32G lane速率，5 mW/lane功耗。
除特别说明，物理参数取最坏情况（TDP + 稳态），不引入瞬态或降额假设。

% 可扩展性
\section{可扩展性：可行性 LP 的规模与求解时间}\label{sec:scalability}

\textbf{论证目标。}
证明框架在 DSE 实用规模内高效可解：可行性 LP（式~\eqref{eq:unified-lp}）的规模随端口数 $N$ 增长受控，
求解时间允许批量 DSE 扫描。这是所有后续实验的前提——若单点求解不可承受，$B^*$ 分析与 Pareto 扫描都无从谈起。

\textbf{实验设计。}
在拓扑家族上枚举一组设计点（Dragonfly 的 $(a,p,h)$ 网格、Mesh、Torus、$k$-ary $n$-cube），
对每个设计点构造并求解完整可行性 LP：性能侧用式~\eqref{eq:valiant-lp} 的 $\min\sum_e L_e$ 压出包络下界，
物理侧含 $\mu$bump 与热 L1。记录三个量：
\begin{enumerate}
    \item 变量数与约束数（LP 规模）随 $N$ 的增长；
    \item 单次 LP 求解时间随规模的增长；
    \item 二分搜索的 LP 调用次数——理论上界 $\log_2(B_{\max}/\varepsilon)$ 与实际次数（第~\ref{sec:ov-bisection}节）。
\end{enumerate}

\textbf{规模构成分析。}
LP 变量以分流变量 $\{\mathbf{f}^{(r)}\}_{r\in\mathcal{R}}$ 为主，规模由排列代表元数 $|\mathcal{R}|$
（附录 A）与端口数 $N$ 共同决定：流量变量数正比于 $|\mathcal{R}|$ 与每模式端口对数的乘积，
物理侧（$\boldsymbol{\ell},\mathbf{P},\mathbf{N}^{\text{sig}},\mathbf{N}^{\text{pwr}},\mathbf{T},\mathbf{x}$）
随链路数与 die 数线性增长。$|\mathcal{R}|$ 由群论归约从 $N!$ 压到分拆数规模（$p(16)=231$，见附录 A），
是 LP 可解的关键。

\textbf{预期方向。}
\begin{itemize}
    \item 单点求解时间应落在 \todo{求解时间量级}，远快于 cycle-accurate 仿真的分钟级开销~\cite{dally2004principles}；
    \item 变量数增长应落在 \todo{规模复杂度} 区间内，$N$ 增至实用规模时不出现超线性爆炸；
    \item 二分搜索的 LP 调用次数应接近 $\log_2(B_{\max}/\varepsilon)$，总开销为单点开销的 \todo{倍数} 倍；
    \item 若某拓扑族的 $|\mathcal{R}|$ 或路径集过大导致 LP 不可解，则回退到附录 A 的更强归约，
    或改用第~\ref{sec:bg-nonblocking}节的 RNB 下界作为可行域下界。
\end{itemize}

% B_max 分析
\section{实验二：真正的物理瓶颈是什么？}\label{sec:bmax}

\textbf{论证目标}：展示框架的区分力——即使拓扑性能相同（$t^* = 1.0$），物理天花板在不同设计点上有数量级的差异。

\textbf{论证逻辑}：
如果只问"800 Gbps下是否feasible"，所有Dragonfly配置返回相同的答案。
但这是初筛工具的失败——它应该能区分设计点的优劣。
因此我们换一个问法：\textbf{不是"在800G下是否可行"，而是"最多能撑到多高带宽才被物理约束杀死"。}

\textbf{实验设计}：
对9个拓扑求解Valiant LP，取最优$\mathbf{L}^*$，代入Eq.~\ref{eq:bmax}计算$B_{\max}^{\text{geom}}$和$B_{\max}^{\text{thermal}}$。
绘制$B_{\max}$--$N$曲线，标注先绑定约束。

\textbf{预期发现}：
\begin{enumerate}
    \item \textbf{Bump预算始终先触顶}。$B_{\max}^{\text{geom}} < B_{\max}^{\text{thermal}}$在所有9个拓扑上成立，
    散热余量2--7倍。这意味着信号bump先用完——优先投资bump pitch缩减而非更激进的冷却。
    \item \textbf{$B_{\max} \propto 1/N$}。从$N=2$的391K Gbps到$N=24$的67K Gbps。
    可外推：$N=96$时$B_{\max} \sim 17$K Gbps。
    \item \textbf{在$N \le 30$和800 Gbps目标下，物理约束有80倍以上余量}。
    架构师可以安全地先专注于拓扑优化，物理验证后置。
\end{enumerate}

\textbf{这个实验的核心价值}：它展示了统一LP框架的区分力——不是"这个设计可不可行"，而是"这个设计的极限在哪里"。
这一信息直接指导技术路线规划。

\begin{figure}[!htbp]
    \centering
    \includegraphics[width=0.85\textwidth]{bmax.pdf}
    \caption{$B_{\max}$随$N$的变化。实线为bump预算上限$B_{\max}^{\text{geom}}$，虚线为散热上限$B_{\max}^{\text{thermal}}$。阴影区域为散热余量。}
    \label{fig:bmax}
\end{figure}

\begin{table}[!htbp]
    \centering
    \caption{B$_{\max}$分析——9个拓扑的物理带宽天花板}
    \label{tab:bmax}
    \footnotesize
    \begin{tabular}{lrrrl}
        \hline
        拓扑 & $N$ & $B_{\max}^{\text{geom}}$ (Gbps) & $B_{\max}^{\text{thermal}}$ (Gbps) & 先绑定约束 \\
        \hline
        DF(1,1,1) & 2  & 391,544 & 1,830,400 & geometry \\
        DF(2,1,1) & 6  & 161,453 & 468,523  & geometry \\
        DF(2,2,1) & 12 & 112,766 & 306,923  & geometry \\
        DF(2,3,1) & 18 & 87,718  & 229,895  & geometry \\
        DF(3,2,1) & 24 & 66,779  & 133,916  & geometry \\
        DF(2,4,1) & 24 & 71,389  & 183,405  & geometry \\
        Mesh(4)    & 16 & 66,925  & 469,295  & geometry \\
        Torus(4)   & 16 & 97,886  & 686,400  & geometry \\
        \hline
    \end{tabular}
\end{table}

% DSE 扫描与 Pareto 前沿
\section{实验三：设计空间长什么样？}\label{sec:pareto}

\textbf{论证目标}：展示框架在批量DSE扫描中的实用价值——在多个设计点上运行LP，揭示Pareto前沿结构，给出设计指导。

\textbf{实验设计}：
固定端口带宽$B = 800$ Gbps，在Dragonfly参数空间$(a,p,h)$中枚举8个配置（$a \le 3, p \le 5$），
每个配置求解Eq.~\ref{eq:valiant-lp} + 几何L1 + 功耗L1。
在$N$--面积平面上绘制Pareto前沿，标注支配关系和总功耗。

\textbf{预期发现}：
\begin{itemize}
    \item Pareto前沿上只有两个支配点：DF(1,1,1)（$N=2$, 288mm$^2$）和DF(2,5,1)（$N=30$, 432mm$^2$）
    \item 在固定面积下，增加每路由器终端数$p$是提升$N$的最有效策略——
    DF(2,2,1)到DF(2,5,1)，面积不变，$N$从12到30
    \item 增加$h$或$a$增加裸片数量和总面积，效率低于增大$p$
\end{itemize}

\textbf{设计指导}：优先增大$p$，再考虑增大$a$或$h$。

\begin{figure}[!htbp]
    \centering
    \includegraphics[width=0.8\textwidth]{pareto.pdf}
    \caption{Dragonfly设计空间的Pareto前沿（$N$ vs 面积，$B=800$Gbps，Valiant路由）。星标为Pareto前沿点，气泡大小表示总功耗。}
    \label{fig:pareto}
\end{figure}

% 精度充足性
\section{实验四：精度充足性——什么时候L0就够了？}\label{sec:fidelity}

\textbf{论证目标}：验证梯度配置的实用价值——对于哪些设计点，低精度已经足够？对于哪些，需要升级？

\textbf{论证逻辑}：
梯度配置的价值不在于"我们有多个精度"，而在于"不同场景用不同精度的判断是一致的，
且在已知安全裕度的情况下可以自信地使用低精度"。

\textbf{实验设计}：
在$N=2, 6, 12, 18, 24$的Dragonfly配置上，分别用$(0,0,0)$和$(1,1,1)$精度求解，
比较$t^*$和feasibility判定。核心指标：L0的$t^*$相比L1是更保守（更安全）还是更乐观（有风险）？

\textbf{预期发现}：
\begin{itemize}
    \item L0始终更保守（更高的$t^*$，更低的$B_{\max}$）——二分带宽下的下界比Valiant分流更松
    \item 这意味着L0给出的"不可行"判断是可靠的（通过L0淘汰的，L1一定也淘汰），
    但L0的"可行"判断需要L1验证（L0可行的，L1可能仍不可行）
    \item 当L0的slack大于某个阈值时（如100×），升级到L1不改变判定，L0精度充足
\end{itemize}

\textbf{框架的使用模式}：先跑L0 $\to$ 观察slack $\to$ 仅当slack低于阈值时升级到L1。
精度升级成为数据驱动的决策，而非方法论绑定。

